% =====================================================================
%  RESUMEN EJECUTIVO EN INGLÉS  ---  1 PÁGINA / 200 A 400 PALABRAS
% =====================================================================
%  Cubre los cinco criterios de la categoría A de la rúbrica:
%    [1] contextualiza empresa, área y puesto      -> párrafo 1
%    [2] objetivo de la empresa para la práctica   -> párrafo 2
%    [3] actividades, metodologías y resultados    -> párrafos 3 y 4
%    [4] inglés correcto, voz activa, tiempo pasado
%    [5] síntesis de contexto, objetivos, metodología, resultados
%        y conclusiones                            -> párrafo 5
%
%  OJO CON LA EXTENSIÓN: tiene que caber en UNA plana. Si agregas texto,
%  vuelve a compilar y revisa que la Tabla de Contenido siga empezando en la
%  plana 3 del PDF.
% =====================================================================

\selectlanguage{British}

\noindent{\Large\textbf{Abstract}\par\vspace{1ex}}

This report describes the professional internship carried out between 2 January
and 28 February 2026 at Desarrollo de Tecnologías y Sistemas SpA (DTS), a
Chilean electronics company serving the defence, security and aerospace markets.
The internship took place in the Research and Development department, as an
Electrical Engineering intern, on a space development line that the
company's strategic plan defined but had never executed.

DTS pursued two objectives during this period: taking its first formal step
towards space capabilities, and making use of the technology transfer agreement
signed with INVAP. With no staff available to open a new line of work, the
company started the project with one staff engineer and one intern. Both objectives converged on one
deliverable, namely bringing the payload design up to Preliminary Design Review
maturity before committing resources to detailed design.

BIFROST covers the architectural design of a Store \& Forward communications
payload for a 12U CubeSat, which receives autonomous
messages from remote ground stations in UHF, stores them on board and downloads
them over Chilean territory. The work followed a sequential systems engineering
life cycle structured around formal design reviews, with IEEE Std 15288.2-2014
defining the deliverables and exit criteria of each gate. The intern modelled the
system with the ARCADIA method on Eclipse Capella and developed a custom orbital
and link budget simulator in Python.

The internship closed a baseline of 13 user and 21 system requirements, linked
by 35 traceability relations with full bidirectional coverage. Twelve of them
were rated as high risk and received a mitigation measure, and the intern
specified 19 verification tests. The
simulation confirmed that the link closes above the 14 dB Eb/N0 threshold on
four to five daily passes, with contact windows of about ten minutes. It also
revealed a zenith null in the ground antenna pattern that constrains download
planning and that a constant-gain calculation would have missed.

The project closed the System Requirements Review and the Preliminary Design
Review within the period, confirming the architecture as technically feasible. DTS kept a reusable system model, a simulation tool and a
documented methodology, started training its engineers in space projects, and
continues developing BIFROST towards the Critical Design Review.

\selectlanguage{Spanish}
\newpage
